\documentclass[10pt,a4paper]{amsart}
\usepackage[margin=19mm]{geometry}
\usepackage{amsmath,amssymb,mathtools}
\usepackage[scr=rsfs,cal=euler]{mathalfa}
\usepackage{xfrac}
\usepackage[unicode,hidelinks,bookmarks=false]{hyperref}
\newcommand{\ie}{i.e.\ }
\newcommand{\cf}{cf.\ }
\newcommand{\resp}{respectively\ }
\newcommand{\Subsection}{\subsection}
\providecommand{\nobreakdash}{\nobreak\hbox{-}}
\setlength{\emergencystretch}{2em}
\multlinegap0pt
\usepackage{tikz}
\usetikzlibrary{cd}
\usepackage{amssymb}
\usepackage{mathtools}
\usepackage[shortlabels]{enumitem}
\usepackage{xcolor}
\usepackage{tikz}
\newtheorem{theorem}{Theorem}
\newtheorem{proposition}{Proposition}
\newcommand{\Aut}{\operatorname{Aut}}
\newcommand{\LG}{\operatorname{LG}}
\DeclareMathOperator{\PSp}{PSp}
\newcommand{\ev}{\operatorname{ev}}
\DeclareMathOperator{\unb}{unb}
\title{Compactifications of spaces of symmetric matrices and pointed Kontsevich spaces of isotropic Grassmannians}
\author{Hanlong Fang, Alex Massarenti, Xian Wu}
\date{Source: arXiv:2603.05808v1 (CC BY 4.0)}
\begin{document}
\maketitle

\noindent\emph{Source and licence.} Compactifications of spaces of symmetric matrices and pointed Kontsevich spaces of isotropic Grassmannians. Version arXiv:2603.05808v1 (CC BY 4.0), distributed by arXiv under the Creative Commons Attribution 4.0 International licence, http://creativecommons.org/licenses/by/4.0/, as recorded in the arXiv licence metadata for this exact version and shown on the abstract page https://arxiv.org/abs/2603.05808v1. Full licence notice: \texttt{licenses/arxiv-2603.05808-CC-BY-4.0.txt}.\par\smallskip

\noindent\textit{Editorial scope.} Counted unit: the article's proposition prop:AutKontsevichConicsLGpointed (source0.tex lines 3005--3010). The article's complete original proof of it is delivered with it (source0.tex lines 3012--3059, one \texttt{proof} environment of 48 lines). No mathematics is written, repaired or re-derived: the statement and the proof are the article's own lines, in the article's own notation, and no retained line is substituted or re-flowed.  Retained uncounted as the source-local context the proof needs: lines 2746--2751; lines 2777--2784; lines 2976--2981.  Nothing was omitted from any retained span, so the package carries no omission record.  No retained line contains a citation, so the wrapper carries no bibliography and no bibliography entry.  The retained lines contain the article's own diagram code, which the wrapper typesets with the standard tikz packages; no image file is delivered and the package declares no asset dependency.  Editorial formatting only, in addition: an AMS \texttt{amsart} preamble that restates the article's own declarations and macros used by the retained lines, namely the environments \texttt{theorem}, \texttt{proposition} on separate counters, together with the standard packages those lines need.  The retained environments print in this document as Theorem 1, Proposition 1 in order of appearance, whereas the article prints the same items with the numbering of its own full document.  The complete native source of the article as distributed in the arXiv e-print for this exact version is bundled unchanged as \texttt{sources/195823/source0.tex}.
\begin{theorem}\label{thm:Eff-general-n-k1}
For every $n\ge 2$, the cone of effective divisors on $\overline{M}_{0,1}(\LG(n,2n),2)$ is given by
\[
\operatorname{Eff}(\overline{M}_{0,1}(\LG(n,2n),2))\ =\ \bigl\langle H_1, \Delta_{1|1}, D_{\unb}\bigr\rangle.
\]
\end{theorem}

\begin{theorem}\label{thm:nef-pointed-conics-LG}
For every $n\ge 2$ the nef cone of $\overline{M}_{0,1}(\LG(n,2n),2)$ is given by
\[
\operatorname{Nef}(\overline{M}_{0,1}(\LG(n,2n),2))
=
\langle H_1, H_{\sigma_2}, T\rangle.
\]
\end{theorem}

\begin{proposition}\label{prop:AutKontsevichConicsLG}
Automorphisms of $\overline{M}_{0,0}\big(\LG(n,2n),2\big)$ are induced by that of the target $\LG(n,2n)$:
\[
\Aut(\overline{M}_{0,0}(\LG(n,2n),2))\cong\PSp_{2n},\,\,\forall\, n\geq2.
\]
\end{proposition}

\begin{proposition}\label{prop:AutKontsevichConicsLGpointed}
Automorphisms of $\overline{M}_{0,1}(\LG(n,2n),2)$ are induced by that of $\LG(n,2n)$:
\[
\Aut(\overline{M}_{0,1}(\LG(n,2n),2))\cong\PSp_{2n},\,\,\forall\,n\geq 2.
\]
\end{proposition}

\begin{proof}
By Theorems \ref{thm:Eff-general-n-k1}, \ref{thm:nef-pointed-conics-LG}, any
$\varphi\in\Aut(\overline{M}_{0,1}(\LG(n,2n),2))$ preserves the extremal rays of these cones. In particular, $\varphi^{*}$ preserves the ray
$\mathbb{R}_{\ge 0}\langle H_1\rangle$, hence the complete linear system $|mH_1|$ for $m\gg 0$.
Because $H_1=\ev_1^{*}H$, the morphism defined by $|mH_1|$ is exactly $\ev_1$, so there is a unique
$\widehat\varphi\in \Aut(\LG(n,2n))$ such that $\ev_1\circ \varphi=\widehat\varphi\circ \ev_1$.
Moreover, $H_{\sigma_2}$ and $T$ on $\overline{M}_{0,1}(\LG(n,2n),2)$ are pullbacks via the forgetful morphism
$\phi_1$ from $\overline{M}_{0,0}(\LG(n,2n),2)$, so preserving the corresponding contraction yields a unique
$\bar\varphi\in \Aut(\overline{M}_{0,0}(\LG(n,2n),2))$ with $\phi_1\circ \varphi=\bar\varphi\circ \phi_1$.  Therefore, $\varphi$ induces an automorphism $\bar\varphi\in\Aut(\overline{M}_{0,0}(\LG(n,2n),2))$.


Identifying $\bar\varphi$ with an element of $\PSp_{2n}$ via Proposition~\ref{prop:AutKontsevichConicsLG}, we derive a natural homomorphism
\[
\chi\colon\Aut(\overline{M}_{0,1}(\LG(n,2n),2))\rightarrow \PSp_{2n}
\qquad
\chi(\varphi):=\bar\varphi.
\]
Moreover, 
$\varphi\in\Aut(\overline{M}_{0,1}(\LG(n,2n),2))$ and  $\bar\varphi=\chi(\varphi)$ fit into the following commutative diagrams.
\[
\begin{tikzcd}
\overline{M}_{0,1}(\LG(n,2n),2) \arrow[r,"\varphi"] \arrow[d,"\phi_1"'] &
\overline{M}_{0,1}(\LG(n,2n),2) \arrow[d,"\phi_1"] \\
\overline{M}_{0,0}(\LG(n,2n),2) \arrow[r,"\bar\varphi"'] &
\overline{M}_{0,0}(\LG(n,2n),2)
\end{tikzcd}
\,\, {\rm and}\,\,
\begin{tikzcd}
\overline{M}_{0,1}(\LG(n,2n),2) \arrow[r,"\varphi"] \arrow[d,"\ev_1"'] &
\overline{M}_{0,1}(\LG(n,2n),2) \arrow[d,"\ev_1"] \\
\LG(n,2n) \arrow[r,"\bar\varphi"'] &
\LG(n,2n).
\end{tikzcd}
\]
Here $\phi_1$ and $\ev_1$ are the forgetful and evaluation morphisms, respectively.


Now let $\varphi\in\ker(\chi)$. Then for every point $[f,x]\in \overline{M}_{0,1}(\LG(n,2n),2)$, we have
\[
\phi_1(\varphi([f,x]))=\phi_1([f,x])=[f],
\qquad
\ev_1(\varphi([f,x]))=\ev_1([f,x])=f(x).
\]
That is,  $\varphi([f,x])$ belongs to the intersection
\(
\pi^{-1}([f]) \cap \ev_1^{-1}(f(x))
\), which consists of the unique point $[f,x]$. Consequently $\varphi([f,x])=[f,x]$ for every $[f,x]$, and thus $\chi$ is an isomorphism.
\end{proof}

\end{document}
